% =============================================================================
\section{维度 2：路由器层}
\label{sec:router}
% =============================================================================

路由器位于客户端与推理资源池之间。我们的论文覆盖了三类自适应机制：静态规则、在线 bandit 自适应，以及基于强化学习的选择。

\subsection{静态语义路由}
\label{sec:router:static}

\vsr{} 立场论文~\cite{vllmsr2026}建立了由信号驱动的架构：把十三类信号组合成适配具体部署的策略，并覆盖 vLLM、OpenAI、Anthropic、Azure、Bedrock 与 Gemini 等后端。嵌入相似度信号由 mmBERT-Embed~\cite{mmberted2025}提供，它的二维 Matryoshka 设计允许运营者按部署需要调节延迟与质量（第 6 层时为 2.56\,ms，适合实时路由；第 22 层时为 11\,ms，适合最大召回）。ProbPol~\cite{probpol2026}则形式化了这些信号的冲突检测问题。

令牌预算资源池路由~\cite{poolrouting2026}实现了一种专门的静态规则：根据估算的总令牌预算，把请求分发到短池或长池，从而减少 17--39\% 的 GPU。OATS~\cite{oats2026}针对工具选择问题，通过零额外服务成本的嵌入优化，将 NDCG@5 从 0.869 提升到 0.940。

WhenToReason~\cite{wang2025reason}另外负责判断某个查询是否值得进入 reasoning 模式，从而避免简单问题上不必要的 chain-of-thought token。Factcheck Classifier~\cite{factcheck2026}又添加了一种预路由信号：被判断为事实查询的请求，可以被导向幻觉率更低的模型，或者在响应后交给 HaluGate~\cite{halugate2025}做进一步验证。

对于多轮会话，\vsr{} 项目~\cite{vllmsr2026}除了单请求分类以外，还提供了两种机制。第一种是\emph{上下文记忆注入}：路由器会累积前几轮中的上下文信号，例如主题、领域、用户画像与难度变化轨迹，并在后续轮次的路由决策中注入这些上下文，从而让推理端点拿到更丰富的背景并提升回答准确率。这一点在精神上类似于 OpenAI Responses API 这类多轮 API 中的会话状态管理。第二种是\emph{面向多轮安全的对比嵌入分类}：路由器不是孤立地判别每一轮，而是利用 few-shot 对比嵌入来检测跨多轮分散投放对抗内容的越狱尝试，而这类攻击是单轮分类器无法发现的。

内容安全与隐私又构成了另一道预路由门控。MLCommons Safety Level-1 分类器~\cite{safetyl1_2026}对每个入站请求进行二元安全/不安全检查；被标记的请求会传递给 Level-2 危害分类器~\cite{safetyl2_2026}，后者从九个 MLCommons 类别中选出其一（暴力犯罪、非暴力犯罪、性犯罪、武器/CBRNE、自残、仇恨、专业建议、隐私、错误信息）。这种分层设计让安全流量的额外延迟几乎可以忽略不计（只有不安全请求才支付第二阶段开销），同时也使按类别进行路由成为可能，例如把侵犯隐私的提示在网关直接阻断，把错误信息查询导向事实依据更强的模型，并把专业建议请求记录下来用于合规审计。这两个分类器都使用 mmBERT LoRA 适配器，符合本项目对低资源占用的设计约束。

FastRouter~\cite{fastrouter2026}保证所有这些路由决策，无论是单轮还是多轮，都能在延迟预算内完成：端到端 50\,ms，GPU 占用 800\,MB，小到足以与服务一起共置。

\subsection{面向聊天的在线自适应路由}
\label{sec:router:online}

聊天工作负载天然带有一种反馈信号：用户的下一条消息会表明上一条回复是否令人满意。我们的 mmBERT-32K Feedback Detector~\cite{feedbackdet2026}把每个用户轮次分成四类（满意、需要澄清、答案错误、想换一种方式），准确率达到 98.8\%，额外开销低于 1\,ms。当检测器判断用户不满意时，路由器可以把后续轮次改派给更强的模型或不同的资源池；这里不需要显式奖励模型，因为用户自身的反应就构成了奖励。再结合路由器的上下文记忆注入机制（\Cref{sec:router:static}），这种自适应就变成了会话感知的：路由器不只是知道用户不满意，还知道是哪个主题和哪个难度层级触发了不满意，因此可以执行有针对性的重路由，而不是一刀切地升级模型。

在我们的项目之外，RouteLLM~\cite{ong2025routellm}使用人类偏好数据训练路由器（以 26\% 的成本达到 95\% 的 GPT-4 质量）；MixLLM~\cite{lu2025mixllm}使用带查询标签的上下文 bandit（以 24.18\% 的成本达到 97.25\% 的 GPT-4 质量）；SageServe~\cite{jia2025sageserve}则把流量预测与基于 ILP 的路由结合起来（在微软规模下降低 25\% 的 GPU 小时）。

\subsection{面向智能体的 RL 路由}
\label{sec:router:rl}

智能体工作负载不像聊天那样拥有逐轮用户反馈：“用户”往往是另一个 LLM 或编排器，而任务成功只有在多步会话结束时才能观测到。这样的延迟且稀疏的奖励，使 RL 比 bandit 更合适。Router-R1~\cite{routerr1_2025}展示了这一思路：它把路由器本身实例化为一个 LLM，在内部推理的 “think” 动作与模型调用的 “route” 动作之间交替，并用 PPO 端到端训练，奖励函数同时包含格式、结果与成本，从而使路由器无需手工启发式就能学到性能-成本权衡。R2-Router~\cite{r2router2026}联合选择模型和输出长度预算，把成本降低到原来的 1/4 到 1/5。M-CMAB~\cite{mcmab2026}则在异构预算下，对多模态调度使用在线 RL。

我们的 AVR 系统~\cite{avr2026}把级联机制应用到了 VLM 智能体任务：第 1 层是多模态分类器（约 $\sim$10\,ms），第 2 层是带 token 过滤评分的小型 VLM 置信探测，第 3 层是大模型升级。它预计可把成本降低到多达 78\%，同时与全大模型基线相比仍保持在 2 个百分点以内。当它与 Visual Confused Deputy 防护~\cite{vcd2026}结合时，高风险动作会直接升级到最强模型。

\paragraph{自适应机制取决于工作负载。}
对\emph{聊天}来说，由 Feedback Detector~\cite{feedbackdet2026}这类分类器检测到的逐轮用户反馈，提供了一种快速、低开销的在线调整信号。对\emph{智能体}来说，只有在多步会话结束后才能看到任务级结果，因此更适合使用基于 RL 的优化~\cite{routerr1_2025}，以对一系列路由决策进行推理。

\subsection{智能体循环、记忆与集群级编排}
\label{sec:router:agent-loops}

近期的智能体服务研究指出，面向工具使用型智能体的推理服务应被视为一个\emph{数据系统}问题，而不是一串互不相关的聊天补全。Helium~\cite{helium2026}围绕工作流结构、复用和算子放置重构了智能体 LLM 服务；Sutradhara~\cite{sutradhara2026}对工具密集型推理中的编排器和执行引擎做联合设计；CONCUR~\cite{concur2026}把拥塞感知批处理用于高吞吐智能体推理；AgServe~\cite{agserve2025}表明\emph{会话感知}调度能突破静态成本-质量权衡；Continuum~\cite{continuum2025}则把多轮调度和 KV-cache 的生存时间耦合起来，避免被暂停的工具调用悄悄破坏局部性。这些系统与基于经验的路由（EvoRoute~\cite{evoroute2026}）以及预算感知顺序路由（BAAR~\cite{baar2026}）是对齐的，但它们通常都只在单一栈内部进行评估。

语义路由器对它们形成补充，因为它导出了一个\emph{全局集群级}控制面：同一套信号驱动的模型路由钩子，可以同时做到：(i) 按 ITR 风格~\cite{itr2026}重塑工具与指令载荷；(ii) 附加记忆层决策，例如全量上下文、摘要、或以检索优先的 prefill，这一过程受 MemCost~\cite{memcost2026} 与长程智能体上的 provider prompt-caching 研究~\cite{dontbreakcache2026}启发；(iii) 发出调度提示，如延后接纳、批处理亲和，从而减少由排队诱发的重试，而后者常常是额外智能体轮次的隐藏来源。xRouter 的 RL 编排~\cite{xrouter2025}说明了，只要奖励定义在完整轨迹之上，路由层就能学到成本感知策略；如果再把这种能力与 OpenClaw~\cite{openclaw2026}这类网关原生栈结合，其中真实用户通过单一控制平面驱动高方差的多轮流量，那么可用于学习的轨迹多样性就会成倍增加，而无需所有框架共用一个记忆系统或工具 SDK。
